% Biuletyn stanu upraw (B1 + B2) — szablon LaTeX wypełniany przez B1_B2/bulletin/build_bulletin.py (Jinja2).
% Znaczniki Jinja2: VAR = wstawiona wartość, BLOCK = logika (pętla, warunek) — zamiast podwójnych klamr, bo klamry są w LaTeX.
% Wszystkie teksty z danych są „oczyszczone” z %, &, _, #, $ w build_bulletin.py (filtr tex).
\documentclass[10pt,a4paper,twocolumn]{article}
\usepackage[T1]{fontenc}
\usepackage[utf8]{inputenc}
\usepackage[polish]{babel}
\usepackage{lmodern}
\usepackage[a4paper,top=1.3cm,bottom=1.5cm,left=1.4cm,right=1.4cm,columnsep=0.7cm]{geometry}
\usepackage{graphicx}
\usepackage{booktabs}
\usepackage{xcolor}
\usepackage{tabularx}
\usepackage{titlesec}
\usepackage{fancyhdr}
\usepackage{lastpage}
\usepackage{microtype}
\usepackage{caption}

\definecolor{worse}{HTML}{D7301F}
\definecolor{normal}{HTML}{2E8B57}
\definecolor{better}{HTML}{2C6FBB}
\definecolor{late}{HTML}{E6862E}
\definecolor{muted}{HTML}{6B6B6B}
\definecolor{boxbg}{HTML}{F2F4F3}

\titleformat{\section}{\large\bfseries}{\thesection.}{0.5em}{}
\titlespacing*{\section}{0pt}{1.0ex}{0.5ex}
\setlength{\parindent}{0pt}
\setlength{\parskip}{0.45ex}

\pagestyle{fancy}
\fancyhf{}
\renewcommand{\headrulewidth}{0pt}
\fancyfoot[L]{\footnotesize\color{muted} Biuletyn stanu upraw · \VAR{county_name|tex} · wydanie \VAR{edition|tex}}
\fancyfoot[R]{\footnotesize\color{muted} strona \thepage{} z \pageref{LastPage}}

\newcommand{\kpi}[3]{%
  \colorbox{boxbg}{\begin{minipage}[t][1.55cm]{0.225\textwidth}\centering
    \vspace{0.15cm}{\footnotesize\color{muted}#1}\\[0.1cm]
    {\Large\bfseries #2}\\[0.05cm]{\scriptsize #3}\end{minipage}}}

\begin{document}

\twocolumn[{%
  {\Large\bfseries Biuletyn stanu upraw}\hfill{\small\color{muted} wydanie \VAR{edition|tex}}\\[0.2cm]
  {\large \VAR{county_name|tex}} {\color{muted}· województwo \VAR{voiv_name|tex} · TERYT \VAR{region_id|tex}}\\[0.1cm]
  {\small Termin: \textbf{\VAR{period_label|tex}} \quad {\color{muted}(zdjęcia satelitarne MODIS, kompozyt 16-dniowy)}}\\[0.25cm]
  \kpi{Zieloność pól}{\VAR{pct_of_norm|tex}}{\VAR{pct_caption|tex}}\hfill
  \kpi{Stan upraw}{\textcolor{\VAR{status_color}}{\VAR{status_short|tex}}}{\VAR{status_caption|tex}}\hfill
  \kpi{Start sezonu}{\VAR{start_kpi|tex}}{\VAR{start_caption|tex}}\hfill
  \kpi{Faza sezonu}{\VAR{phase_short|tex}}{\VAR{phase_caption|tex}}
  \vspace{0.35cm}
}]

\section*{Streszczenie}
\VAR{summary|tex}

{\footnotesize\color{muted}\VAR{summary_note|tex}}

\section{Stan upraw (B1)}
\VAR{b1_text|tex}

\begin{center}\footnotesize
\begin{tabularx}{\columnwidth}{@{}Xrrl@{}}
\toprule
Termin & \% normy & $z$ & Stan \\
\midrule
\BLOCK{for r in b1_rows}
\VAR{r.label|tex} & \VAR{r.pct|tex} & \VAR{r.z|tex} & \textcolor{\VAR{r.color}}{\VAR{r.status|tex}} \\
\BLOCK{endfor}
\bottomrule
\end{tabularx}
\end{center}
{\footnotesize\color{muted} Norma: \VAR{norm_years|tex}, bez lat skrajnych. $z$ = odchylenie od normy w jednostkach typowego rozrzutu; $|z|\ge 1$ oznacza wyraźne odchylenie.}

\begin{figure}[h]
  \centering
  \includegraphics[width=\columnwidth]{fig_curve.pdf}
  \caption*{\footnotesize Krzywa NDVI pól ornych w \VAR{year} r. na tle normy (pasmo: 25--75\% typowych lat). Pionowe linie: start, szczyt i koniec sezonu (B2); kreska przerywana: termin wydania.}
\end{figure}

\section{Przebieg sezonu (B2)}
\VAR{b2_text|tex}

\begin{center}\footnotesize
\begin{tabularx}{\columnwidth}{@{}Xllr@{}}
\toprule
Etap & Data & Typowo & Przesunięcie \\
\midrule
\BLOCK{for r in b2_rows}
\VAR{r.stage|tex} & \VAR{r.date|tex} & \VAR{r.typical|tex} & \textcolor{\VAR{r.color}}{\VAR{r.shift|tex}} \\
\BLOCK{endfor}
\bottomrule
\end{tabularx}
\end{center}
\BLOCK{if group_rows}
{\footnotesize Grupy upraw (osobne krzywe z pikseli jednej uprawy, rok mapy upraw \VAR{year}):}
\begin{center}\footnotesize
\begin{tabularx}{\columnwidth}{@{}Xlll@{}}
\toprule
Grupa & Start & Szczyt & Koniec \\
\midrule
\BLOCK{for r in group_rows}
\VAR{r.group|tex} & \VAR{r.start|tex} & \VAR{r.peak|tex} & \VAR{r.end|tex} \\
\BLOCK{endfor}
\bottomrule
\end{tabularx}
\end{center}
\BLOCK{endif}

\begin{figure}[h]
  \centering
  \includegraphics[width=\columnwidth]{fig_map.pdf}
  \caption*{\footnotesize Stan upraw w powiatach województwa (\VAR{voiv_name|tex}) w tym terminie (B1). \VAR{neighbours_text|tex}}
\end{figure}

\section{Struktura upraw}
\VAR{structure_text|tex}

\begin{center}
  \includegraphics[width=0.95\columnwidth]{fig_structure.pdf}
\end{center}

\section{Warunki pogodowe}
\VAR{weather_text|tex}

\section{Ocena i możliwe przyczyny}
\VAR{assessment_text|tex}

\section*{Metoda i źródła}
{\footnotesize
NDVI: MODIS MOD13Q1 (NASA, 250\,m, co 16 dni), piksele pól ornych wg mapy upraw EUCROPMAP (JRC, \VAR{map_year}), bez chmur i śniegu.
Krzywa wygładzona (Savitzky--Golay), luki do 48 dni uzupełniane. Norma: \VAR{norm_years|tex}, lata dalsze niż 1,5 odchylenia odrzucane.
Fenologia: start i koniec sezonu = przejście przez 50\% wysokości rocznej „górki” NDVI, szczyt = maksimum (dokładność ok. $\pm$8 dni).
Granice: GUGiK (PRG). Struktura upraw: EUCROPMAP (nie rozróżnia ozimin i jarych — grupy przybliżone).
Ograniczenia: piksel 250\,m często łączy kilka pól; wynik dotyczy powiatu, nie pojedynczego gospodarstwa.
\VAR{ai_note|tex}
Wygenerowano automatycznie: \VAR{generated|tex}.\par}

\end{document}
